\title{Controlling Text-to-Image Diffusion by Orthogonal Finetuning}

\begin{abstract}
State-of-the-art text-to-image diffusion models demonstrate remarkable effectiveness in synthesizing photorealistic images from textual descriptions. A central challenge lies in the systematic guidance and adaptation of these sophisticated models to accommodate a variety of downstream tasks. To this end, we present a theoretically grounded adaptation approach—Orthogonal Finetuning~(OFT)—designed for refining text-to-image diffusion models for diverse tasks. Distinct from prior strategies, OFT offers a theoretical guarantee for conserving hyperspherical energy, which encapsulates the inter-neuronal relationships on the unit hypersphere. We demonstrate that maintaining this property is essential to preserving the semantic integrity of generative outputs. For increased stability during model adaptation, we introduce Constrained Orthogonal Finetuning (COFT), which augments OFT by enforcing a radius restriction on the hypersphere. Our investigation focuses on two prominent text-to-image finetuning scenarios: subject-driven generation, which entails synthesizing images of a particular subject from a few sample images alongside a textual prompt; and controllable generation, where the model incorporates supplementary control signals during synthesis. Experimental evaluation indicates that the OFT approach achieves superior generation fidelity and accelerated convergence compared to current techniques.
\end{abstract}

\section{Introduction}

Emerging text-to-image diffusion frameworks~\cite{saharia2022photorealistic,ramesh2022hierarchical,rombach2022high} have set new standards in producing high-resolution images conditioned on textual input. Despite these advancements, relying solely on text for conditioning often introduces ambiguity, limiting the extent of precise, fine-level manipulation in generated outputs. This paper addresses two specific categories within text-to-image generation that reflect this limitation:

\begin{itemize}[leftmargin=*,nosep]

    \item \textbf{Subject-driven generation}~\cite{ruiz2023dreambooth}: Given a small set of sample images depicting a target subject, the goal is to synthesize novel images representing the same subject in alternate scenarios as guided by a text description.
    \item \textbf{Controllable generation}~\cite{zhang2023adding,mou2023t2i}: Provided with external control information (such as canny edges or segmentation masks), this task focuses on producing images that both conform to the specified control and adhere to a text-based prompt.
\end{itemize}

The central challenge in both tasks revolves around how to adapt text-to-image diffusion models through finetuning while retaining the generative quality established during pretraining. For a finetuning method to be considered effective, it should fulfill the following criteria: (1) \emph{training efficiency}, which involves reducing both the number of trainable parameters and the required training epochs, and (2) \emph{generalizability preservation}, maintaining the model's ability to generate diverse, high-fidelity outputs. Finetuning approaches generally fall into two categories: directly adjusting the model’s weights using a small learning rate (\emph{e.g.}, \cite{ruiz2023dreambooth}), or integrating a compact module with its own re-parameterized weights (\emph{e.g.}, \cite{hulora2022,zhang2023adding}). However, these straightforward strategies do not guarantee the conservation of generative performance acquired during pretraining. Furthermore, there is an absence of clear principles to inform the design of effective finetuning techniques and the selection of optimal hyperparameters, such as the number of training epochs. One principal challenge stems from the inability to rigorously measure the extent to which pretrained generative capabilities are preserved. Current finetuning approaches are predicated on the assumption that a minimal Euclidean distance between the finetuned and original models corresponds to superior retention of pretrained capabilities, leading practitioners to use small learning rates. Nonetheless, although this may sometimes be adequate, we contend that merely considering Euclidean proximity to the pretrained model fails to comprehensively represent semantic preservation. Employing a more structurally informed metric to capture discrepancies between the finetuned and pretrained models could substantially improve both the safeguarding of original generative performance and the overall stability during finetuning.

Building on empirical findings that hyperspherical similarity effectively represents semantic relationships~\cite{liu2017deep,liu2018decoupled,chen2020angular}, our approach utilizes hyperspherical energy~\cite{liu2018learning} to describe the pairwise relational configuration among neurons within a network layer. Hyperspherical energy, computed as the aggregate of hyperspherical similarities (\emph{e.g.}, cosine similarities) across all neuron pairs within the same layer, serves to reflect the uniformity of neuron distribution over the unit hypersphere~\cite{liu2021learning}. We postulate that an optimal finetuned model should exhibit only minimal deviations in hyperspherical energy relative to its pretrained counterpart. Although one could simply incorporate a regularization term to enforce stability of hyperspherical energy during finetuning, this strategy does not inherently guarantee minimization of energy discrepancies. Consequently, we exploit a key invariance property: applying a uniform orthogonal transformation to all neurons leaves their pairwise hyperspherical similarities unchanged. Based on this property, we introduce Orthogonal Finetuning~(OFT), a method for adapting large-scale text-to-image diffusion models to new tasks while maintaining constant hyperspherical energy. The main mechanism involves learning an orthogonal transformation shared across each layer so that the pairwise angular relationships among neurons are retained. This process can be interpreted as redefining the reference coordinate frame for all neurons in a given layer. Furthermore, noting that maintaining a smaller Euclidean distance between finetuned and pretrained models tends to preserve pretrained capabilities, we extend our method to Constrained Orthogonal Finetuning~(COFT). This variant constrains the finetuned weights to reside within a hypersphere of predetermined radius centered at the pretrained neuron positions.

The rationale underlying the effectiveness of orthogonal transformation for neuron finetuning is partly inspired by the 2D Fourier transform, which separates an image into its magnitude and phase spectra. Notably, the phase spectrum, representing the angular relationship between the input and the basis, encodes the bulk of the image's semantic content. For instance, reconstructing an image using its phase spectrum combined with a randomized magnitude spectrum can still retain the essential semantics of the original. This observation indicates that manipulating neuron orientations is central to semantic alteration in generated imagery—a fundamental objective in both subject-driven and controllable generation tasks. However, if neuron directions are altered with excessive freedom, this can compromise the generative abilities acquired during pretraining. To mitigate this, we propose to constrain these modifications by preserving the pairwise angles between neurons, grounded in the assumption that such angles are pivotal for capturing a neural network’s learned knowledge. Accordingly, we advocate for learning a shared orthogonal transformation per layer, maintaining invariant hyperspherical energy across transformations.

Our approach also takes cues from orthogonal over-parameterized training~\cite{liu2021orthogonal}, which develops classification networks from scratch via orthogonal transformation applied to a randomly initialized architecture. This method leverages the fact that random initializations yield low hyperspherical energy on average, and~\cite{liu2021orthogonal} aims to maintain this low energy throughout training (a factor correlated with enhanced generalization in classification~\cite{liu2018learning,lin2020regularizing}). Their findings demonstrate the sufficiency of orthogonal transformations for training generalizable classifiers. In contrast, our focus is on finetuning text-to-image diffusion models to improve controllability and achieve superior generative outcomes. OFT departs from~\cite{liu2021orthogonal} in two fundamental ways. First, whereas~\cite{liu2021orthogonal} targets minimizing hyperspherical energy, OFT is designed to preserve the hyperspherical energy level characteristic of the pretrained model, safeguarding the network’s internal structure from being disrupted during finetuning. In diffusion model finetuning, minimizing this energy could erase core semantic structures. Second, OFT is constructed to minimize divergence from the pretrained weights, resulting in a constrained optimization variant;~\cite{liu2021orthogonal} does not implement this restriction. Successfully finetuning a model thus depends on a careful equilibrium between adaptability and retention of prior knowledge—an objective that our OFT method is well-positioned to accomplish. Our primary contributions can be summarized as follows:

\begin{itemize}[leftmargin=*,nosep]

    \item We introduce a new approach to finetuning, termed Orthogonal Finetuning (OFT), designed to enhance the controllability of text-to-image diffusion models. To further bolster training stability, we develop a constrained version of the method that restricts the angular displacement from the pretrained parameters.
    \item In contrast to standard finetuning strategies, OFT uniquely maintains the hyperspherical energy of the model parameters during training, a property we theoretically establish and empirically show to be crucial for retaining the pretrained model’s generative semantics.
    \item We validate OFT on two prominent use cases: subject-driven and controllable image generation. Through a thorough experimental analysis, we show that our method markedly outperforms previous techniques in terms of output quality, speed of convergence, and finetuning robustness. OFT also demonstrates greater data efficiency, requiring fewer training samples and epochs to achieve convergence.
    \item For evaluating controllable generation, we present a novel metric called control consistency, which quantifies how well the generated image reflects the intended control signal by inferring and comparing the signal from the output with the original input. Our results highlight OFT’s superior controllability as confirmed by this metric.
\end{itemize}

\section{Related Work}

\textbf{Text-to-image diffusion models}. Considerable advances~\cite{saharia2022photorealistic,ramesh2022hierarchical,rombach2022high,gu2022vector,nichol2022glide} have been achieved in text-conditioned image generation, largely due to significant developments in diffusion-based generative approaches~\cite{ho2020denoising,song2021score,song2021denoising,dhariwal2021diffusion} and multimodal vision-language representation learning~\cite{su2020vlbert,lu2019vilbert,radford2021learning,wang2021simvlm,li2022blip,li2023blip,alayrac2022flamingo,schuhmann2022laion}. Systems like GLIDE~\cite{nichol2022glide} and Imagen~\cite{saharia2022photorealistic} perform training in pixel space. GLIDE incorporates a joint optimization of the text encoder with a diffusion prior leveraging paired text-image datasets, whereas Imagen freezes a pretrained text encoder during diffusion model training. In contrast, models such as Stable Diffusion~\cite{rombach2022high} and DALL-E2~\cite{ramesh2022hierarchical} learn in a latent feature space. Stable Diffusion utilizes VQ-GAN~\cite{esser2021taming} to obtain a compressed visual codebook as its latent space, while DALL-E2 leverages CLIP~\cite{radford2021learning} to construct a shared image-text latent embedding space. Apart from diffusion-based frameworks, alternative generative paradigms—such as generative adversarial networks~\cite{reed2016generative,zhang2017stackgan,xu2018attngan,li2019controllable} and autoregressive models~\cite{ramesh2021zero,ding2021cogview,wu2022nuwa,yuscaling}—have also been investigated for text-to-image synthesis. OFT is formulated as a model-agnostic finetuning strategy, applicable to a broad range of text-to-image diffusion models.

\textbf{Subject-driven generation}. To address the challenge of unwanted subject alteration, some works~\cite{avrahami2022blended,nichol2022glide} propose incorporating user-provided masks as supplementary conditioning signals. Inversion-based techniques~\cite{choi2021ilvr,dhariwal2021diffusion,ramesh2022hierarchical,gal2022image} have been developed to enable context editing without modifying the main subject, and other approaches~\cite{hertz2022prompt} can perform both localized and holistic image edits absent explicit mask input. Nevertheless, these solutions may struggle to retain fine identity features of the depicted subject. Pivotal Tuning~\cite{roich2022pivotal} involves finetuning a generator around an initially inverted latent code, applying regularization to help preserve subject identity. A related strategy~\cite{nitzan2022mystyle} learns individualized generative priors using a person’s facial image set. Casanova et al.~\cite{casanova2021instance} present a method to generate variations of a specific instance, though at times instance-specific characteristics may be lost. DreamBooth~\cite{ruiz2023dreambooth} enhances a text-to-image diffusion model using a few images and a dedicated token, combining reconstruction with class-specific

\textbf{Controllable generation}. Image-to-image translation can be regarded as a particular instance of controllable generation, with most established approaches relying on conditional generative adversarial networks~\cite{isola2017image,zhu2017toward,wang2018high,park2019semantic,choi2018stargan}. More recently, diffusion models have been introduced as alternative solutions for image-to-image translation tasks~\cite{saharia2022palette,voynov2022sketch,wang2022pretraining}. Advances such as ControlNet~\cite{zhang2023adding} offer a method for incorporating control signals into a pretrained diffusion model through finetuning, achieving highly effective controllable image synthesis. Likewise, T2I-Adapter~\cite{mou2023t2i}, developed concurrently, also leverages finetuning strategies for pretrained diffusion models to strengthen the controllability of outputs. In line with the methodologies of \cite{zhang2023adding,mou2023t2i}, we employ OFT for finetuning diffusion models, and consistently observe enhanced controllability even when using limited training samples and fewer tunable parameters. Notably, OFT maintains efficiency by incurring no extra inference-time computational costs.

\textbf{Model finetuning}. Adapting large pretrained models to suit downstream tasks through finetuning has gained notable traction in the research community~\cite{devlin2018bert,brown2020language,he2022masked}. Adaptation techniques, such as those introduced in \cite{hulora2022,houlsby2019parameter,pfeiffer2020adapterfusion}, have been extensively investigated in natural language processing as special cases of finetuning. OFT shares a close relationship with LoRA~\cite{hulora2022}, which employs a low-rank paradigm for modeling weight updates. In contrast, OFT implements a multiplicative update of neuron weights via layer-wise orthogonal transformations, which analytically guarantee the preservation of pairwise angular relationships between neurons within the same layer and thus contribute to improved stability.

\section{Orthogonal Finetuning}

\subsection{Why Does Orthogonal Transformation Make Sense?}

We begin our exploration by motivating the role of orthogonal transformation in the finetuning process of text-to-image diffusion models. Our analysis focuses on two fundamental aspects: (1) the rationale for finetuning the angular (directional) components of neuron weights, and (2) the justification for employing orthogonal transformations as the mechanism for this angular adjustment.

To address the first question, we build on empirical findings reported in \cite{liu2018decoupled,chen2020angular}, which highlight that angular discrepancies between features effectively represent the semantic gap. SphereNet~\cite{liu2017deep} provides evidence that constraining all neurons of a neural network to reside on a unit hypersphere preserves expressivity and can even enhance generalization, suggesting that the directional component of neuron representations encapsulates the most critical data information. To further illustrate the significance of neuron angles, we design a simplified experiment, depicted in Figure~\ref{fig:toy_ae}. In this experiment, a typical convolutional autoencoder is trained on a dataset of flower images. During training, the conventional inner product is utilized to compute the feature map ($z$ represents the output of the convolutional kernel $\bm{w}$ when applied to input $\bm{x}$ in the moving window). For the testing phase, three methods are examined for producing the feature map: (a) the same inner product as in training, (b) retaining only magnitude information, and (c) utilizing only angular information. The visualizations in Figure~\ref{fig:toy_ae} reveal that neuron angular information alone is almost sufficient to reconstruct the input images, whereas neuron magnitude does not contribute meaningful information. It is crucial to note that the cosine activation (c) is not introduced during training; all training leverages only the inner product. These findings indicate that the primary vehicle for encoding the semantic content of input images is the directionality (angles) of neurons. Thus, adjustments to neuron directions are expected to more efficiently influence or modify image semantics.

In response to the second question, a direct approach for adapting neuron directions involves applying a collective rotation or reflection to all neurons within a given layer, which inherently invokes an orthogonal transformation. While it is theoretically possible to employ more intricate angular transformations—such as rotating individual neurons by varying amounts—we observe that orthogonal transformations offer an optimal compromise between adaptability and structural regularity. Supporting this, \cite{liu2021orthogonal} demonstrates the substantial representational power of orthogonal transformations for neural network learning. To substantiate the effectiveness of this regularization, we conduct an experiment using the ControlNet~\cite{zhang2023adding} framework to generate controllable outputs, with results summarized in Figure~\ref{fig:ortho}. It is clear from these outcomes that the standard OFT exhibits robust performance, achieving precise control upon the completion of training (epoch 20), while removing the orthogonality constraint from OFT leads to a complete breakdown in generating realistic images or exerting any control. These experimental observations underscore the critical role played by the orthogonality constraint in ensuring the efficacy of OFT.

Traditional finetuning approaches often rely on gradient descent with a reduced learning rate, updating either the entire model or select layers. Employing a small learning rate implicitly constrains updates to remain close to the initial pretrained weights, effectively minimizing $\thickmuskip=2mu \medmuskip=2mu \| \bm{M} - \bm{M}^0\|$, where $\bm{M}$ represents the adapted model’s parameters and $\bm{M}^0$ denotes the original pretrained parameters. While this inherent limitation provides an intuitive safeguard against drastic changes, it may still not sufficiently restrict the search space when dealing with large-scale models. To address this, LoRA introduces a more explicit limitation, imposing a low-rank structure on the weight difference: $\thickmuskip=2mu \medmuskip=2mu \text{rank}(\bm{M} - \bm{M}^0)=r'$, where $r'$ is chosen to be small. In contrast, OFT takes a different approach by regulating the similarity among neuron pairs through a hyperspherical energy constraint: $\thickmuskip=2mu \medmuskip=2mu \| \text{HE}(\bm{M})-\text{HE}(\bm{M}^0) \|=0$, where $\text{HE}(\cdot)$ indicates the hyperspherical energy corresponding to the weight matrix. 

As a concrete instance, take a fully connected layer expressed as $\thickmuskip=2mu \medmuskip=2mu \bm{W}=\{\bm{w}_1,\cdots,\bm{w}_n\}\in\mathbb{R}^{d\times n}$, with $\bm{w}_i\in\mathbb{R}^{d}$ as the $i$-th column vector, representing an individual neuron; $\bm{W}^0$ denotes its pretrained counterpart. The output $\thickmuskip=2mu \medmuskip=2mu \bm{z}\in\mathbb{R}^n$ is computed as $\thickmuskip=2mu \medmuskip=2mu \bm{z}=\bm{W}^\top\bm{x}$, with $\thickmuskip=2mu \medmuskip=2mu \bm{x}\in\mathbb{R}^d$ as input. The OFT method can thus be framed as the minimization of the hyperspherical energy gap between the finetuned and original weights:

\begin{equation}\label{eq:oft_principle}
\footnotesize
\min_{\bm{W}} \left\| \text{HE}(\bm{W})-\text{HE}(\bm{W}^0)\right\|~~\Leftrightarrow~~\min_{\bm{W}} \bigg{\|} \sum_{i\neq j} \|\hat{\bm{w}}_i-\hat{\bm{w}}_j\|^{-1}-\sum_{i\neq j} \|\hat{\bm{w}}^0_i-\hat{\bm{w}}^0_j\|^{-1}\bigg{\|}
\end{equation}

In this equation, $\thickmuskip=2mu \medmuskip=2mu \hat{\bm{w}}_i:=\bm{w}_i/\|\bm{w}_i\|$ defines the normalized vector for neuron $i$, and the layer’s hyperspherical energy is given as $\thickmuskip=2mu \medmuskip=2mu \text{HE}(\bm{W}):= \sum_{i\neq j} \|\hat{\bm{w}}_i-\hat{\bm{w}}_j\|^{-1}$. It is straightforward to see that the objective in Eq.~\eqref{eq:oft_principle} achieves a minimum value of zero, which is attained when $\bm{W}$ results from applying a rotation or reflection to $\bm{W}^0$, i.e., $\thickmuskip=2mu \medmuskip=2mu \bm{W}=\bm{R}\bm{W}^0$ with $\thickmuskip=2mu \medmuskip=2mu \bm{R}\in\mathbb{R}^{d\times d}$ orthogonal (where a determinant of $1$ or $-1$ indicates rotation or reflection, respectively). The core principle behind OFT, therefore, is to constrain model adaptation such that neurons within each layer are transformed by a shared orthogonal

\begin{equation}\label{eq:oft_general}
    \footnotesize
    \bm{z}=\bm{W}^\top\bm{x}=(\bm{R}\cdot\bm{W}^0)^\top\bm{x},~~~\text{s.t.}~\bm{R}^\top\bm{R}=\bm{R}\bm{R}^\top=\bm{I}
\end{equation}

Here, $\bm{W}$ represents the OFT-finetuned parameter matrix, while $\bm{I}$ denotes the identity matrix. An overview of OFT is provided in Figure~\ref{fig:block}. Following the approach of zero initialization used in LoRA, it is critical that OFT adjusts the pretrained model starting precisely from $\bm{W}^0$. For this purpose, the orthogonal matrix $\bm{R}$ is initially set to the identity matrix, which ensures that the updated model parameters coincide with the original pretrained weights at initialization. To maintain the orthogonality constraint on $\bm{R}$ throughout training, we can employ differentiable orthogonalization approaches introduced in \cite{liu2021orthogonal,lezcano2019cheap}. We further elaborate on efficient strategies to preserve orthogonality below.

\subsection{Efficient Orthogonal Parameterization}\label{sect:eff_ortho}

While common orthogonalization procedures, such as the Gram-Schmidt algorithm, are differentiable, they are computationally intensive for large-scale applications~\cite{liu2021orthogonal}. To improve computational efficiency, we utilize the Cayley parameterization method to construct the orthogonal matrix. In this formulation, the orthogonal matrix is given by $\thickmuskip=2mu \medmuskip=2mu \bm{R}=(\bm{I}+\bm{Q})(I-\bm{Q})^{-1}$, where $\bm{Q}$ is a skew-symmetric matrix, i.e., $\thickmuskip=2mu \medmuskip=2mu \bm{Q}=-\bm{Q}^\top$. Although the Cayley parameterization is computationally advantageous, it is limited to generating orthogonal matrices with determinant $1$, which are elements of the special orthogonal group. However, empirical results indicate that this restriction does not hinder practical performance. Nevertheless, for large values of $d$, the full parameterization of $\bm{R}$ using the Cayley transform can still be inefficient. To mitigate this issue, we propose structuring $\bm{R}$ as a block-diagonal matrix with $r$ independent blocks, yielding:

\begin{equation}
    \footnotesize
    \bm{R}=\text{diag}(\bm{R}_1,\bm{R}_2,\cdots,\bm{R}_r)=\begin{bmatrix}
    \bm{R}_1\in \text{O}(\frac{d}{r}) &  &\\
    &\ddots & \\
    & & \bm{R}_r\in \text{O}(\frac{d}{r})
    \end{bmatrix}\in \text{O}(d)
\end{equation}

Here, $\text{O}(d)$ refers to the set of $d$-dimensional orthogonal matrices. Specifically, $\thickmuskip=2mu \medmuskip=2mu \bm{R}\in\mathbb{R}^{d\times d}$ as well as each $\thickmuskip=2mu \medmuskip=2mu \bm{R}_i\in\mathbb{R}^{d/r\times d/r},\forall i$, are orthogonal by construction. Setting $\thickmuskip=2mu \medmuskip=2mu r=1$ recovers the standard, unconstrained orthogonal matrix form as the block-diagonal structure reduces to a full matrix. For a general orthogonal matrix of size $\thickmuskip=2mu \medmuskip=2mu d\times d$, there are $\thickmuskip=2mu \medmuskip=2mu d(d-1)/2$ independent parameters, corresponding to a parameter complexity of $\mathcal{O}(d^2)$. For an orthogonal matrix composed of $r$ blocks, each block of size $d/r\times d/r$, the total parameter count becomes $\thickmuskip=2mu \medmuskip=2mu d(d/r-1)/2$, giving a parameter complexity of $\mathcal{O}(d^2/r)$. Moreover, further reductions in parameterization can be obtained by sharing the block matrix across all blocks, i.e., enforcing $\thickmuskip=2mu \medmuskip=2mu \bm{R}_i = \bm{R}_j, \forall i\neq j$, which in turn reduces the parameter complexity down to $\mathcal{O}(d^2/r^2)$. Throughout these approaches for parameter savings, the orthogonality of $\bm{R}$ is retained, thereby ensuring that the underlying hyperspherical energy property is unaffected.

A comparison of parameter efficiency between OFT and LoRA is warranted. In LoRA, for a given low-rank parameter $r'$, the trainable parameter count amounts to $\thickmuskip=2mu \medmuskip=2mu r'(d+n)$. When $r$ and $r'$ are proportional to the hidden dimension $d$ (e.g., $\thickmuskip=2mu \medmuskip=2mu r=r'=\alpha d$, for some $0<\alpha\leq 1$), LoRA's parameter complexity is $\mathcal{O}(d^2 + dn)$, while for OFT it becomes $\mathcal{O}(d)$. To concretely demonstrate this contrast, consider a weight matrix of size $\thickmuskip=2mu \medmuskip=2mu 128\times 128$. With $\thickmuskip=2mu \medmuskip=2mu r'=8$, LoRA introduces $2,048$ trainable parameters, whereas OFT, also with $\thickmuskip=2mu \medmuskip=2mu r=8$ and no block sharing, utilizes $960$ trainable parameters.

\subsection{Constrained Orthogonal Finetuning}

The expressiveness of OFT can be deliberately reduced by restricting the finetuned model to stay within a predefined vicinity of the original pretrained solution. This leads to the Constrained Orthogonal Finetuning (COFT) scheme, whose computation is defined as:

\begin{equation}\label{eq:coft_general}
    \footnotesize
    \bm{z}=\bm{W}^\top\bm{x}=(\bm{R}\cdot\bm{W}^0)^\top\bm{x},~~~\text{s.t.}~\bm{R}^\top\bm{R}=\bm{R}\bm{R}^\top=\bm{I},~\norm{\bm{R}-\bm{I}}\leq \epsilon
\end{equation}

In this formulation, the model is subject to both an orthogonality constraint and a restriction enforcing $\bm{R}$ to stay within an $\epsilon$-distance of the identity matrix. While the Cayley parameterization from Section~\ref{sect:eff_ortho} ensures orthogonality, embedding the additional $\epsilon$-deviation constraint within this parameterization is challenging. To analyze the structure of the Cayley transform, we use the Neumann series expansion to approximate $\thickmuskip=2mu \medmuskip=2mu \bm{R}=(\bm{I}+\bm{Q})(I-\bm

Because the series converges with respect to the operator norm, the restriction $\thickmuskip=2mu \medmuskip=2mu\|\bm{R}-\bm{I}\|\leq\epsilon$ can be incorporated directly into the Cayley transform framework. Under this reparameterization, the equivalent requirement becomes $\thickmuskip=2mu \medmuskip=2mu \|\bm{Q}-\bm{0}\|\leq\epsilon'$, where $\bm{0}$ is the zero matrix and $\epsilon'$ is a new error tolerance (potentially distinct from $\epsilon$). This modified constraint on $\bm{Q}$ is straightforward to enforce via projected gradient descent methods. To ensure the orthogonal matrix $\bm{R}$ starts as the identity, we initialize $\bm{Q}$ to be the zero matrix. The COFT approach, therefore, integrates two explicit constraints: minimal change in hyperspherical energy, and a controlled deviation from the pretrained model parameters. While existing finetuning techniques often leverage the second constraint implicitly, COFT formalizes it explicitly. Empirically, although OFT demonstrates strong performance, COFT offers enhanced stability during finetuning, attributed to the explicit deviation constraint. Figure~\ref{fig:eps_coft} shows the impact of the choice of $\epsilon$ on COFT performance. The results indicate that $\epsilon$ mediates the trade-off in adaptation: larger $\epsilon$ values make COFT behavior closely resemble that of OFT, while smaller $\epsilon$ make the resulting model act increasingly like the original pretrained text-to-image diffusion system.

\subsection{Re-scaled Orthogonal Finetuning}

We introduce a straightforward augmentation to standard OFT by also learning a scaling factor for each individual neuron. This enhancement draws motivation from the observation that scaling neuron activations does not affect hyperspherical energy, since normalization eliminates scaling effects. Concretely, our forward computation becomes $\thickmuskip=2mu \medmuskip=2mu\bm{z}=(\bm{R}\bm{W}^0\bm{D})^\top\bm{x}$\footnote{\textbf{Errata}: The NeurIPS camera-ready version inaccurately states the forward computation as $\thickmuskip=2mu \medmuskip=2mu\bm{z}=(\bm{D}\bm{R}\bm{W}^0)^\top\bm{x}$. However, the codebase correctly implements the intended operation; as a result, findings reported in Appendix~\ref{app:rescaled} remain valid.} where the scaling matrix $\thickmuskip=2mu \medmuskip=2mu\bm{D}=\text{diag}(s_1,\cdots,s_n)\in\mathbb{R}^{n\times n}$ is a diagonal matrix whose diagonal entries $s_1,\cdots,s_n$ are learnable and strictly positive. Unlike the original OFT formulation in Eq.~\eqref{eq:oft_general} which restricts adaptation to the matrix $\bm{R}$ alone, this formulation allows both $\bm{R}$ and $\bm{D}$ to be optimized. This modification brings about increased expressiveness at the cost of very few extra parameters. Nevertheless, we adhere to the original OFT for our main experiments to highlight the standalone merit of orthogonal transformations, although we note that in practice, the re-scaled OFT generally yields stronger results (refer to Appendix~\ref{app:rescaled} for more detail).

\section{Intriguing Insights and Discussions}

\textbf{OFT accommodates a variety of model architectures}. In theory, OFT is universally applicable to diverse neural network architectures. For example, while LoRA is conventionally employed on Transformer attention weights~\cite{hulora2022}, to facilitate a fair comparison in our experiments, we limit OFT to finetuning attention matrices. Moreover, OFT is particularly amenable to convolutional layers due to the unique interpretability of the block-diagonal structure of $\bm{R}$ in this context—an advantage over LoRA. When configuring the number of blocks in $\bm{R}$ to match the input channel count, each block exclusively updates an individual neuron channel, mirroring the effect of depth-wise convolution~\cite{chollet2017xception}. Alternatively, when the same orthogonal block is shared among all channels, OFT transforms all channels uniformly using a common orthogonal matrix. Preliminary experiments on OFT for convolutional layer adaptation can be found in Appendix~\ref{app:convolution}.

\textbf{Connection to LoRA}. LoRA limits drastic alterations to the pretrained weight matrix by introducing a low-rank update, ensuring that most of the original information is preserved. In comparison, OFT enforces orthogonality (with full rank) on the transformation applied to the pretrained weight matrix, thereby safeguarding the integrity of the pretraining information by not distorting it. The OFT forward computation can be reformulated as $\thickmuskip=2mu \medmuskip=2mu \bm{z}=(\bm{R}\bm{W}^0)^\top\bm{x}=(\bm{W}^0+(\bm{R}-\bm{I})\bm{W}^0)^\top\bm{x}$, where $\thickmuskip=2mu \medmuskip=2mu (\bm{R}-\bm{I})\bm{W}^0$ serves a comparable role to LoRA's low-rank adaptation. Given that $\bm{W}^0$ is generally full-rank, OFT effectively implements a low-rank modification whenever $\thickmuskip=2mu \medmuskip=2mu \bm{R}-\bm{I}$ is constrained to be low-rank. Much like LoRA uses a rank parameter $r'$, OFT utilizes a block-diagonal parameter $r$ to restrict the number of trainable parameters. Importantly, LoRA and OFT exemplify two fundamentally different strategies for achieving parameter efficiency: LoRA harnesses a low-rank parameterization to curtail the parameter count, whereas OFT capitalizes on structural sparsity—in particular, block-diagonal orthogonality.

\textbf{Why OFT converges faster}. Firstly, as demonstrated in Figure~\ref{fig:toy_ae}, modifying neuron orientations is the most direct approach for inducing semantic changes, which is precisely the function OFT fulfills. Additionally, OFT may be interpreted as adapting neurons within the smooth geometry of a hypersphere manifold, a property that provides a more favorable optimization landscape. This observation is further corroborated by empirical results in \cite{liu2021orthogonal}.

\textbf{Why not minimize hyperspherical energy}. Our approach diverges from \cite{liu2021orthogonal} in that we do not attempt to optimize hyperspherical energy. In typical classification tasks, the objective is to avoid redundant neurons. While minimizing hyperspherical energy would distribute neurons evenly across the hypersphere, this uniformity does not constitute a desirable goal in the context of finetuning, as it can overwrite or erode essential information acquired during pretraining.

\textbf{Balancing flexibility and regularization in finetuning}. Our findings highlight a fundamental trade-off between flexibility and regularity. While conventional finetuning provides maximum adaptability, it is prone to instability and often leads to model degradation. Remarkably, the OFT approach, despite its straightforward design, strikes an effective compromise between adaptability and constraint by maintaining the angular relationships between neurons. Furthermore, the block-diagonal scheme serves to enforce an even greater level of regularization on the orthogonal matrix.

\textbf{Inference efficiency maintained}. In contrast to ControlNet, the OFT method does not introduce any extra computational burden during inference. At test time, the learned orthogonal matrix $\bm{R}$ can be directly post-multiplied with the original pretrained weights $\bm{W}^0$, resulting in a transformed weight matrix $\thickmuskip=2mu \medmuskip=2mu \bm{W}=\bm{R}\bm{W}^0$. Consequently, the inference runtime remains unchanged from that of the original model.

\section{Experiments and Results}

\textbf{Experimental configuration}. The experiments are carried out using Stable Diffusion v1.5~\cite{rombach2022high} as the base text-to-image generation model. For unbiased assessment, images generated by each method are selected at random. The experimental protocols for subject-driven tasks are adapted from DreamBooth~\cite{ruiz2023dreambooth}, while the setup for controllable image synthesis is based on ControlNet~\cite{zhang2023adding} and T2I-Adapter~\cite{mou2023t2i}. To provide an equitable comparison with LoRA, both OFT and COFT are restricted to modifying only those layers to which LoRA is applied. Additional experimental results and further methodological details can be found in Appendix~\ref{app:settings}.

\subsection{Subject-driven Generation}

\textbf{Experimental setup}. For benchmarks, we use DreamBooth~\cite{ruiz2023dreambooth} and LoRA~\cite{hulora2022}. All approaches are optimized using the same loss employed in DreamBooth. In the cases of DreamBooth and LoRA, we generally adhere to the procedures and hyperparameter choices recommended in their respective original publications. Supplementary results and discussions are presented in Appendix~\ref{app:settings},\ref{app:coft_oft},\ref{app:more_qualitative},\ref{app:failure}.

\textbf{Finetuning stability and convergence}. To begin, we assess the stability and convergence rate during finetuning for DreamBooth, LoRA, OFT, and COFT, as illustrated in Figure~\ref{fig:teaser} and Figure~\ref{fig:db_convergence}. It is evident that both COFT and OFT are able to finetune diffusion models in a stable manner. By 400 iterations, DreamBooth and OFT-based methods already attain effective control, while LoRA struggles to retain the identity of the subject. Progressing to 2000 iterations, DreamBooth starts exhibiting mode collapse with degraded images, and LoRA is unable to generate a yellow shirt, mistakenly outputting yellow fur instead. Conversely, OFT and COFT continue to maintain stable and consistent regulation of the generated outputs. This demonstrates the ability of our OFT framework to rapidly and reliably converge in the context of subject-driven generation. A noteworthy advantage is the low sensitivity of OFT and COFT to the specific number of finetuning steps, which simplifies hyperparameter selection and obviates extensive iteration tuning for different subjects. Both methods permit the use of a large, fixed iteration count. Moreover, when COFT is equipped with an appropriate $\epsilon$, both the learning rate and the number of iterations can be set with minimal concern.

\textbf{Quantitative comparison}. In line with \cite{ruiz2023dreambooth}, we further provide quantitative assessments covering subject fidelity (using DINO~\cite{caron2021emerging}, CLIP-I~\cite{radford2021learning}), fidelity to text prompts (using CLIP-T~\cite{radford2021learning}), and sample diversity (via LPIPS~\cite{zhang2018unreasonable}). CLIP-I evaluates the mean pairwise cosine similarity of CLIP embeddings between generated samples and authentic images. DINO employs a similar computation, substituting in ViT S/16 DINO embeddings for comparison. For CLIP-T, we calculate the average cosine similarity of CLIP embeddings between the prompt and generated outputs. We also assess sample diversity by calculating the mean LPIPS cosine similarity between outputs corresponding to the same subject and identical prompt. The summary of findings in Table~\ref{table:dreambooth_final} reveals that both COFT and OFT surpass DreamBooth and LoRA by a substantial margin in DINO and CLIP-I, while also exhibiting equal or slightly superior performance for prompt fidelity and sample diversity. For rigorous evaluation, each method is tested via repeated finetuning of the identical pretrained model using 30 independent random seeds to control for randomness. Collectively, these findings confirm that our OFT framework delivers enhanced convergence, greater stability, and more robust overall results.

\textbf{Qualitative comparison}. To provide clearer insight into the advantages offered by OFT, we present several randomly selected cases of subject-driven generation in Figure~\ref{fig:dreambooth_final}. For consistency and fairness, all samples are produced by applying each technique to the same finetuned model, without independently optimizing text prompts for particular outcomes. For every approach, we report the results of the model that attains the highest validation CLIP score. Examination of Figure~\ref{fig:dreambooth_final} shows that both OFT and COFT exhibit strong capabilities for maintaining the semantic features of the subject, whereas LoRA frequently struggles to retain subject identity (\emph{e.g.}, in the bowl instance, LoRA entirely loses the subject). Furthermore, OFT and COFT allow for much more precise alignment with input prompts, in contrast to DreamBooth, which, despite preserving subject identity well, often cannot synthesize outputs adhering to the specified prompt (\emph{e.g.}, the top row of the bowl example). These qualitative results suggest that our OFT approach excels at achieving both effective subject preservation and prompt-based control. Additionally, OFT's performance remains robust regardless of the number of training iterations, unlike DreamBooth and LoRA, both of which are sensitive to this parameter and underperform with more iterations. Expanded qualitative evidence is available in Appendix~\ref{app:more_qualitative}. Human studies presented in Appendix~\ref{app:human} offer further confirmation of OFT’s superior effectiveness.

\subsection{Controllable Generation}

\textbf{Settings}. We compare our method against established baselines ControlNet~\cite{zhang2023adding}, T2I-Adapter~\cite{mou2023t2i}, and LoRA~\cite{hulora2022}. The primary paper investigates three demanding tasks in controllable generation: translating Canny edges to images (C2I) with the COCO dataset~\cite{lin2014microsoft}, generating images from segmentation maps (S2I) on ADE20K~\cite{zhou2017scene}, and reconstructing faces from landmarks (L2F) using CelebA-HQ~\cite{karras2018progressive,xia2021tedigan}. For all approaches, Stable Diffusion~(SD) v1.5 is fine-tuned over 20 epochs on each of these datasets. Additional results can be found in Appendix~\ref{app:more_qualitative}, \ref{app:more_control_task}, and \ref{app:failure}.

\textbf{Convergence}. We analyze the convergence characteristics of ControlNet, T2I-Adapter, LoRA, and COFT in the context of the L2F task, incorporating both numerical and visual assessments. For quantitative comparison, we calculate the mean $\ell_2$ distance between the predicted facial landmarks and the corresponding control landmarks. In Figure~\ref{fig:face_convergence}, we display the model’s face landmark errors as a function of the number of fine-tuning epochs. Our findings indicate that both COFT and OFT demonstrate markedly more rapid convergence than alternative methods. Notably, LoRA requires 20 epochs to attain a level of performance that OFT already achieves at epoch 8. Importantly, OFT and COFT possess a comparable number of trainable parameters to LoRA, which is considerably fewer than those in ControlNet, yet their convergence is substantially more efficient. Additional evidence for the accelerated convergence of OFT is presented in Figure~\ref{fig:teaser}: the example on the right illustrates that OFT surpasses ControlNet and LoRA in data efficiency, successfully converging using merely 5\% of the ADE20K dataset. 

For qualitative analysis, our comparison centers on OFT, COFT, and ControlNet, as ControlNet’s landmark error is most similar to ours. As visualized in Figure~\ref{fig:face_convergence_qual}, both OFT and COFT exhibit consistent and stable convergence, with synthesized facial poses incrementally matching the target landmarks over epochs. In contrast, ControlNet shows a sudden emergence of controllability only after the eighth epoch, mirroring the abrupt convergence dynamics reported by \cite{zhang2023adding}. This enhanced stability in convergence behavior makes the OFT framework substantially more practical, providing smoother and more predictable training dynamics and thereby simplifying the selection of optimal hyperparameters for OFT.

\textbf{Quantitative comparison}. To assess the efficacy of controllable generation, we propose a metric for control consistency. This metric functions by extracting the control signal from the output image and juxtaposing it with the original control input. In the C2I scenario, Intersection-over-Union (IoU) and F1 score are utilized, while the S2I setting employs mean IoU as well as both mean and overall accuracy. For the L2F task, the average $\ell_2$ distance between the control landmarks and those predicted in the generated image is measured. Additional information on these consistency metrics is provided in Appendix~\ref{app:settings}. All competing methods are evaluated using their optimal hyperparameter configurations. Table~\ref{table:control_gen} demonstrates that OFT and COFT achieve notably higher levels of control precision compared to alternative approaches. It is evident that adapter-based strategies such as T2I-Adapter and ControlNet not only converge at a slower rate but also underperform in the final evaluation. LoRA outperforms ControlNet in S2I but falls short in both C2I and L2F. Generally, LoRA's upper-bound performance appears limited despite extensive hyperparameter optimization. In contrast, OFT has yet to reach its maximal capacity, as performance continues to improve with increasing numbers of trainable parameters. It should be highlighted that our approach to quantitatively evaluating controllable generation represents an important contribution, providing precise assessment of control ability for fine-tuned models, subject to the limitations of the chosen off-the-shelf segmentation or detection backbone.

\textbf{Qualitative comparison}. In addition to quantitative evaluations, we conduct a qualitative assessment of OFT and COFT, positioning them against leading methods such as ControlNet, T2I-Adapter, and LoRA. The randomly sampled outputs depicted in Figure~\ref{fig:control_final} illustrate that OFT and COFT not only maintain superior realism and high-fidelity in the generated images, but also offer precise control. For the S2I task, it is evident that LoRA fails to adhere to the input segmentation map, whereas ControlNet, OFT, and COFT all demonstrate robust controllability. Compared with ControlNet, OFT and COFT produce images that are richer in detail and geometric consistency, all while relying on a much smaller number of model parameters. In the C2I task, both OFT and COFT successfully synthesize semantically consistent images from coarse Canny edge inputs, whereas T2I-Adapter and LoRA deliver inferior results. Regarding the L2F task, our approach attains the highest accuracy in pose-conditioned face generation, even with challenging viewpoints. Across all three tasks, our qualitative evaluations indicate that OFT and COFT surpass the state-of-the-art baselines in visual quality, thus confirming the efficacy of the OFT framework for controlled image synthesis. Further diverse qualitative comparisons covering the three control tasks are presented in Appendix~\ref{app:control_exp}. Additionally, Appendix~\ref{app:more_control_task} demonstrates the applicability of OFT to a wider range of control scenarios, such as dense pose to human body, sketch to image, and depth to image translation.

\section{Concluding Remarks and Open Problems}

Inspired by the insight that the angular relationships among neurons play a pivotal role in visual semantics, we introduce a straightforward yet powerful finetuning technique—orthogonal finetuning—for exerting control over text-to-image diffusion models. Our approach specifically addresses two key scenarios: subject-driven generation and controllable generation. When compared to prior approaches, OFT achieves superior control and greater finetuning stability while utilizing fewer finetuned parameters. Crucially, OFT imposes no additional computational burden during inference, making the resulting model highly suitable for practical deployment.

The introduction of OFT also highlights several intriguing open questions. Firstly, OFT enforces orthogonality using Cayley parametrization, which necessitates computing a matrix inverse. This requirement presents a bottleneck for scaling up OFT. While our use of block diagonal parametrization offers a workaround, developing faster and differentiable methods for matrix inversion remains an open challenge. Secondly, OFT offers distinct promise in compositionality, as the orthogonal matrices obtained from several OFT finetuning tasks can be composed through multiplication and retain their orthogonal property. Whether these combined matrices effectively encapsulate the knowledge from all associated downstream tasks is an unresolved issue that warrants investigation. Lastly, the parameter efficiency of OFT is intimately linked to the block diagonal design, which unavoidably injects bias and curtails flexibility. Discovering strategies to boost parameter efficiency with minimal bias and greater effectiveness stands as a significant open problem.